\exnumber{01}
\extitle{UX III - Ex 2 - Questionnaires}
\exfiles{ex01/PDF, nomme Name\_FirstName\_DeliverableName\_Date\_VersionNumber}
\exauthorize{Ceux nommés dans l'exercice}
\makeheaderfiles

\textbf{Instructions :}

Préparez un questionnaire pour collecter des données quantitatives auprès de vos utilisateurs.

Et rappelez-vous, vous n'êtes pas votre utilisateur !

Pour être utilisé dans des études sérieuses et obtenir des résultats viables, un formulaire doit recueillir au moins 100 réponses. Les réponses doivent vous donner des données quantifiables que vous pouvez utiliser pour prouver un point et confirmer des tendances, pas des verbatims ou des détails.

\textbf{Livrables :}

\begin{itemize}
\item Un document PDF comprenant votre ou vos objectif(s) et vos hypothèses, les questions, une introduction et une conclusion.
\item Dans le document, ajouter le lien vers un questionnaire prêt à être diffusé
\item Déposez le PDF sur Github, nommé ainsi : « Name\_FirstName\_DeliverableName\_Date\_VersionNumber ». Exemple : « Doe\_John\_MidFiPrototype\_05242024\_V1 ».
\item Ajoutez un titre dans le document.
\item Soignez l'aspect global du document. Il doit être clair, simple et lisible. Vous pouvez regarder canva.com, sans surcharger le dessin de détails.
\end{itemize}

\textbf{Pistes :}

Voici les étapes recommandées :

1. Définissez votre objectif, ce que vous cherchez en diffusant votre formulaire aux gens. Habituellement, cela concerne des chiffres, des proportions, etc.
2. Notez 1 à 3 hypothèses que vous avez sur ce sujet, que vous devez soit confirmer soit infirmer.
3. Choisissez 3 à 5 sujets sur lesquels il vous faut de l'information.
4. Pour chaque sujet, rédigez 1 à 3 questions fermées. Rappelez-vous, votre utilisateur doit répondre par un choix dans une liste de réponses préparées ou utiliser l'option « autre » pour écrire la réponse qu'il ne trouve pas dans la liste préparée.
5. Rédigez quelques questions de screening pour obtenir quelques données démographiques (âge, métier, localisation, etc.). Cela vous permet de vous assurer que les réponses obtenues viennent de votre cible.
6. Structurez vos questions en entonnoir : commencez par des questions larges pour vous concentrer sur des sujets ou situations spécifiques. N'oubliez pas de rester neutre dans votre ton pour laisser l'utilisateur raconter sa propre histoire. Essayez de vous concentrer sur 10 questions principales.
7. Rédigez une introduction de salutation, décrivant brièvement (1) qui vous êtes et (2) ce que vous cherchez. Pas besoin d'en dire trop, mais assez pour fournir un cadre et l'objectif du formulaire.
8. Rédigez une conclusion pour remercier votre utilisateur pour son temps, et préparez les prochaines étapes. Ex. : demander de diffuser ce formulaire à 3 personnes, etc.

\textbf{Obligatoire : demandez-leur s'ils seraient disponibles pour un entretien anonyme en personne (par téléphone, en présentiel, en visioconférence : selon ce qui les met le plus à l'aise), car vous devrez mener au moins 3 entretiens par la suite.}

\textbf{Ressources :}

\begin{itemize}
\item Lean Survey Canvas : comment l'utiliser \url{https://medium.com/think-big-work-smart/the-lean-survey-canvas-1b0a00cab200}
\item Questionnaires utilisateurs - Comment interroger vos utilisateurs \url{https://refiner.io/blog/user-surveys/}
\item Concevoir un excellent questionnaire en 3 étapes simples \url{https://uxdesign.cc/designing-a-great-survey-with-3-simple-steps-283abaa67ad6}
\item Questionnaires UX : un guide rapide pour en tirer le meilleur \url{https://uxdesign.cc/ux-surveys-a-quick-guide-to-get-the-most-out-of-them-1889d136be6a}
\item Mieux rechercher grâce aux questionnaires \url{https://uxmastery.com/better-user-research-through-surveys/}
\item L'expérience utilisateur des questionnaires \url{https://uxdesign.cc/the-user-experience-of-surveys-a5420c82604a}
\end{itemize}

\textbf{Outils :}

\begin{itemize}
\item Typeform \url{https://www.typeform.com/}
\item Google Forms \url{https://docs.google.com/forms/u/0/}
\end{itemize}

\textbf{Citations :}

\begin{itemize}
\item « Les questionnaires mesurent et catégorisent les attitudes ou collectent des données autodéclarées qui peuvent aider à suivre ou découvrir des problèmes importants à traiter. » Norman Nielsen Group
\item Les questionnaires ne sont pas précis pour fournir des données comportementales car LES UTILISATEURS OMETTENT DES ÉTAPES AU MILIEU ET LES SOUVENIRS SONT DÉFAILLANTS.
\end{itemize}

